\chapter{绪论}

\section{背景意义}
介绍所属研究方向、问题背景及开展综述的意义。

\section{文献概述}
说明文献覆盖范围和分类，概括主要研究分支。上标引文写法见示例\cite{sample-book,sample-article}。

\chapter{研究分支一}

\section{研究现状}
按理论、方法、方案或技术比较代表性工作，说明已有成果解决了哪些问题。

\subsection{代表性方法}
图~\ref{fig:review-framework}为图题及交叉引用的用法示例，可替换为与主题相关的插图。

\begin{figure}[htbp]
  \centering
  \begin{tikzpicture}[box/.style={draw,minimum width=29mm,minimum height=9mm},
    node distance=8mm]
    \node[box] (work) {代表性工作};
    \node[box, right=of work] (result) {已有进展};
    \node[box, right=of result] (gap) {待解难点};
    \draw[-{Latex}] (work) -- (result);
    \draw[-{Latex}] (result) -- (gap);
  \end{tikzpicture}
  \caption{研究现状归纳框架示例}
  \label{fig:review-framework}
\end{figure}

\subsection{方法比较}
表~\ref{tab:review-comparison}示范三线表的写法。

\begin{table}[htbp]
  \centering
  \caption{代表性工作对比示例}
  \label{tab:review-comparison}
  \begin{MasterTable}
  \begin{tabular}{@{}lll@{}}
    \toprule
    对比维度 & 方法一 & 方法二 \\
    \midrule
    适用条件 & 按原文填写 & 按原文填写 \\
    评价依据 & 按原文填写 & 按原文填写 \\
    \bottomrule
  \end{tabular}
  \end{MasterTable}
\end{table}

\section{发展趋势}
讨论制约研究进展的难点及结合已有证据可预见的发展趋势。

\chapter{研究分支二}

\section{研究现状}
按统一维度概括该分支的代表性成果、现有方法的适用条件和待解问题。

\section{发展趋势}
联系不同技术路线的证据与局限，讨论后续研究方向。

\chapter{结论}
综合各研究分支的进展与不足，归纳仍待解决的关键问题，并据此提出本课题拟研究的问题。
